% CEG5206 L02 condensed from chat/ChatGPT-lec2-20261105-2127.md.
% shortcut: original L02.pdf unavailable; verify slide-specific constants when supplied.
\section{L02: Memory Hierarchy \& Cache}
\subsection{Memory, locality \& units}
\textbf{Latency:} time per operation; \textbf{throughput:} achieved work/time;
\textbf{bandwidth:} maximum data/time.
Data throughput $\le$ bandwidth; high bandwidth $\not\Rightarrow$ low latency.
Faster memory is generally smaller and more expensive per byte.
\textbf{Temporal:} same data soon; \textbf{spatial:} nearby data soon.
Reuse after many intervening distinct blocks may miss.
\textbf{CPU:} registers $\to$ L1/L2/L3 $\to$ DRAM; hardware-managed cache.
\textbf{GPU:} per-thread registers; per-block on-chip shared memory (explicit);
hardware L1/L2; global DRAM. Local memory is thread-private DRAM space
(e.g. register spills), not necessarily fast/on-chip.
Constant: broadcast reads; texture: spatial locality.
\textbf{FPGA:} registers, LUTRAM, BRAM/URAM; explicitly designed buffers/dataflow;
external SRAM/DRAM/Flash.
\[
\begin{array}{@{}c@{}}
W(n)=\text{computational work},\\
Q(n)=\text{cache misses},\quad \#\text{lines}=M/\mathcal B.
\end{array}
\]
Cache holds \emph{many} lines; a miss transfers a whole line.
$M$: cache capacity; $\mathcal B$: line size (distinct from matrix $B$).
\textbf{Units:} address mapping uses bytes; matrix analysis uses elements/words.
For $e$ bytes/element, convert both: $M=M_{\rm bytes}/e$,
$\mathcal B=\mathcal B_{\rm bytes}/e$.

\subsection{Cache mapping \& address fields}
Byte address $a$, $w$ address bits, $k$ ways;
$M$ is data capacity (excludes tag/valid metadata).
Assume $\mathcal B$ and set count $E$ are powers of two.
\[
\begin{array}{@{}c@{}}
E=\dfrac{M}{k\mathcal B},\quad
q=\left\lfloor\dfrac a{\mathcal B}\right\rfloor,\\
\text{offset }o=a\bmod\mathcal B,\\
\text{set }s=q\bmod E,\\
\text{tag }t=\lfloor q/E\rfloor=\lfloor a/(M/k)\rfloor,\\
a=[\text{tag}\mid\text{set index}\mid\text{offset}],\\
b_o=\log_2\mathcal B,\quad b_s=\log_2E,\\
b_t=w-b_s-b_o=w-\log_2(M/k).
\end{array}
\]
\textbf{Direct mapped:} $k=1$; unique line/slot;
$b_s=\log_2(M/\mathcal B)$, $b_t=w-\log_2M$.
\textbf{$k$-way:} one fixed set, any of its $k$ ways.
\textbf{Fully associative:} $k=M/\mathcal B$, $E=1$, no index bits;
any line, $b_t=w-\log_2\mathcal B$.
\textbf{Hit:} selected set contains a valid way with matching tag;
offset selects byte, not hit/miss. On miss: use invalid way first,
otherwise evict within the selected set; install block, tag, valid bit.
\textbf{LRU:} hit or fill $\to$ MRU; evict least recently used valid line.
\textbf{Example:} $w=16$, $M=256$ B, $\mathcal B=16$ B, $k=4$:
$E=4$, $(b_t,b_s,b_o)=(10,2,4)$.
$a=\texttt{0x035A}=858$: $q=53$, $(t,s,o)=(13,1,10)$.

\subsection{Miss types \& common traps}
For an actual miss, ignoring coherence:
\textbf{(1) Cold/compulsory:} first access to this \emph{block}, not each element.
\textbf{(2) Capacity:} previously accessed block also misses in a
same-capacity fully associative LRU reference cache.
\textbf{(3) Conflict:} that reference hits, actual mapped cache misses.
Compare the \emph{same trace} and line size; reference policy here is LRU,
not the optimal policy of the ideal-cache model.
\textbf{Traps:} capacity misses can occur in all organizations;
conflict misses occur in direct \emph{and} set associative caches;
fully associative removes mapping conflicts.
A full set/way does not by itself classify the miss; first access stays cold
even if it evicts another block.
\textbf{Sharing/coherence:} another core's write may invalidate a cached line.
True sharing: same data; false sharing: different data on the same line.
Read-only sharing alone does not cause invalidation ping-pong;
padding/separating writable data can reduce false sharing.

\subsection{Double buffering: overlap load/compute}
$N\ge1$ equal chunks; $T_L,T_C>0$ load/compute time per chunk.
Two disjoint buffers alternate: compute current chunk while loading next;
wait for both before swapping/reusing buffers.
Assume independent load/compute resources, serial order within each stage,
and no synchronization/switching overhead (no extra output stage).
\[
\begin{array}{@{}c@{}}
T_{\rm single}=N(T_L+T_C),\\
H=\max(T_L,T_C),\quad L=\min(T_L,T_C),\\
T_{\rm double}=T_L+(N-1)H+T_C\\
=NH+L,\\
\mathrm{Speedup}_N=\dfrac{N(H+L)}{NH+L},\\
\mathrm{Speedup}_{\infty}=1+\dfrac LH\le2.
\end{array}
\]
First load + steady-state overlap + final compute.
$T_L>T_C$: load-bound; $T_C>T_L$: compute-bound.
Balanced stages maximize speedup: if $T_L=T_C$,
$\mathrm{Speedup}_N=2N/(N+1)<2$; approaches $2$ as $N\to\infty$.
$N=1$: no overlap, speedup $1$. More buffers do not remove the slower-stage limit.

\subsection{Ideal cache \& caching lemmas}
\textbf{Ideal:} fully associative, optimal replacement (future knowledge).
\textbf{LRU lemma:} same access trace, line size and initially empty caches:
\[
Q_{\rm LRU}(2M)\le2Q_{\rm Ideal}(M).
\]
LRU has \emph{twice} the capacity; not a same-capacity bound, not equality.
\textbf{Cache-Miss Lemma:} access only $r$ contiguous segments of lengths $S_i$
(segments need not be adjacent), $N=\sum_iS_i$; fully associative cache.
\[
\begin{array}{@{}c@{}}
N<M/3,\quad N/r\ge\mathcal B\\
\Longrightarrow\quad
Q\le N/\mathcal B+2r\le3N/\mathcal B.
\end{array}
\]
Row/segment boundary overhead gives $+2r$; all touched lines fit,
so repeated accesses hit under LRU/ideal. Bound $\ne$ exact count.
\textbf{Tall cache:} $\mathcal B^2<cM$ for a suitably small constant $c>0$;
asymptotically $M=\Omega(\mathcal B^2)$.
\textbf{Submatrix Caching Lemma:} row-major $n\times n$ submatrix;
each row is a contiguous segment ($N=n^2$, $r=n$):
\[
\begin{array}{@{}c@{}}
\mathcal B^2<cM\le n^2<M/3\\
\Longrightarrow\quad Q\le3n^2/\mathcal B.
\end{array}
\]
Lower condition ensures $n>\mathcal B$ (long rows);
upper condition reserves room for row-edge lines.
These are sufficient conditions; smaller submatrices can also cache efficiently.

\subsection{Naive matrix multiplication: three regimes}
\textbf{Assumptions:} row-major $A,B,C$; loops $i\to j\to k$;
fully associative LRU, tall cache, cold start; $M,\mathcal B$ in elements.
\[
C_{ij}\mathrel{+}=A_{ik}B_{kj},\qquad W(n)=\Theta(n^3).
\]
Inner $k$: $A$ contiguous, $B$ stride $n$, $C_{ij}$ reused.
For $n\ge\mathcal B$, one $B$ column touches $n$ distinct lines;
holding these lines needs about $n\mathcal B$ elements, not $n$.
\[
\begin{array}{@{}l@{\quad}l@{}}
\text{Size condition}&Q(n)\\\hline
n>c_1M/\mathcal B&\Theta(n^3)\\
c_2\sqrt M<n<c_3M/\mathcal B&\Theta(n^3/\mathcal B)\\
n<c_4\sqrt M&\Theta(n^2/\mathcal B)
\end{array}
\]
$c_i$ are suitable constants, separately for each regime;
the two threshold \emph{scales} are $\sqrt M$ and $M/\mathcal B$,
not exact numerical cutoffs (transition bands omitted).
Cold-cache expressions above assume $n^2\ge\mathcal B$;
for tiny matrices use $\Theta(1+n^2/\mathcal B)$ in the fitting case.
\textbf{Large:} $n$ current $B$ lines do not fit; little reuse between columns;
fixed $i$: $\Theta(n^2)$ misses, then $n$ values of $i$.
\textbf{Middle:} these $n$ lines fit, entire $B$ does not.
Each group of $\mathcal B$ adjacent columns loads $n$ lines;
fixed $i$: $(n/\mathcal B)n=n^2/\mathcal B$; multiply by $n$ outer iterations.
\textbf{Small:} all three matrices fit; load each line once, reuse across all $i$.
Ignoring row-edge overhead, $3n^2\le M$ is the simple capacity check.
For $n=16$, $\mathcal B=4$, cold $B$ alone needs $256/4=64$ lines,
not $256$; no further factor $n$ if it remains cached.

\subsection{Blocking / tiling: short derivation}
\emph{Supplementary derivation from the lecture's blocking idea.}
$S\times S$ tiles, $\mathcal B\le S\le n$; assume $S\mid n$ for counts,
fully associative cache, and all three actual tile line footprints fit.
Tile rows are contiguous; row-edge overhead is $O(S)$ lines.
\[
\begin{array}{@{}c@{}}
W_{\rm tile}=\Theta(S^3),\quad D_{\rm tile}=\Theta(S^2),\\
Q_{\rm tile}=O(S^2/\mathcal B+S)=O(S^2/\mathcal B),\\
\#\text{tile products}=(n/S)^3,\\
Q_{A+B}=O\!\left((n/S)^3S^2/\mathcal B\right)\\
=O(n^3/(\mathcal BS)).
\end{array}
\]
\textbf{Keep $C$ resident:} outer tile indices $I,J$, inner $K$;
$C_{IJ}\mathrel{+}=A_{IK}B_{KJ}$ for every $K$ before releasing $C_{IJ}$.
\[
\begin{array}{@{}c@{}}
Q_C=O\!\left((n/S)^2S^2/\mathcal B\right)=O(n^2/\mathcal B),\\
Q_{\rm blocked}=O\!\left(\dfrac{n^3}{\mathcal BS}+\dfrac{n^2}{\mathcal B}\right).
\end{array}
\]
Ignoring row-edge overhead: $3S^2\le M$.
For large matrices, choose $S=\Theta(\sqrt M)$ with a small enough constant
to leave line-boundary headroom; tall cache ensures $S\ge\mathcal B$.
\[
Q_{\rm blocked}=O\!\left(\dfrac{n^3}{\mathcal B\sqrt M}+\dfrac{n^2}{\mathcal B}\right).
\]
Work stays $\Theta(n^3)$; reuse per loaded tile grows as $\Theta(S)$.
Too-large tiles evict useful lines; strided access alone does not force a miss
if the needed line is still cached. Write-back adds only a constant factor here.
